%This is a template file for use of iopjournal.cls

\documentclass{iopjournal}

% Options
% 	[anonymous]	Provides output without author names, affiliations or acknowledgments to facilitate double-anonymous peer-review

\begin{document}

\articletype{Paper} %	 e.g. Paper, Letter, Topical Review...

\title{Grid Cost and Emissions Externalities of US Data Center Loads with Temoa}

\author{David J. Florez Rodriguez$^{1}$\orcid{0009-0007-7270-839X}, Paulina Jaramillo$^{1,2}$\orcid{0000-0002-4214-1106} and Joseph Decarolis$^{1}$\orcid{0000-0003-4677-4522}}

\affil{$^1$Engineering and Public Policy, Carnegie Mellon University, Pittsburgh, PA, USA}

\affil{$^2$Civil and Environmental Engineering, Carnegie Mellon University, Pittsburgh, PA, USA}

\affil{$^*$Author to whom any correspondence should be addressed.}

\email{@gmail.com}

\keywords{energy modeling, grid planning, data centers, emissions, }

\begin{abstract}
Data centers may add gigawatts of load to the US grid in the coming years. New generation will be necessary to feed their servers and cooling. The type of generation that grid regulators and industry actors decide to build will impact the nation's grid reliability, public electricity costs, and national emissions. We model possible grid futures exploring `three pillars' regulation that would require data center loads to be met with renewable energy that is (1) additional, (2) hourly-matched, and (3) deliverable. Additional scenarios also explore partial regulation of data centers and uncertainties in future demand, electricity prices, and grid development. These simulations find higher local energy prices and emissions without regulation supporting renewable generation to supply data center loads. If properly regulated, server load growth over future decades can add valuable resiliency to the grid and keep its electricity affordable to all. 
\end{abstract}

\section{Introduction}
Data center loads made up about 1\% of global energy use Cite[masanet,2020] and 2\% of US energy use in 2020 Cite [Shehabi, 2024]. By 2024, these servers consumed almost 5\% of US electricity generation Cite [Shehabi]. The US Energy Information Agency (EIA) estimates they'll consume between 8\% and 10\% of US Electricity in 2050 Cite [EIA 2026]. This large electric load requires large infrastructure investments. Industry reports see 40\% higher annual utility investments in 2026 compared to 2020 Cite [EEI 2025], with server's exponential growth in a time of grid stress as one reason. In the last PJM capacity auction, new server loads have comprised 40\% of the total cost, and the capacity auction fell 6 GW short of their reliability goal Cite [Holand, 2026]. On the ground, data center perception struggles as about \$130B across 75 projects  have been blocked or delayed in 2026 Cite [data center watch, 2026], and dozens of municipalities across various states have passed bans on new data center construction Cite [Route Fifty, 2026]. AI does has some benefits, from Alpha Fold Cite [Abramson, 2024; Nobel Prize, 2024] to improved grid optimization Cite [Baker, 2025; Owerko, 2020]. Nevertheless, since data centers impact the grid, environment, and the lives around these systems, regulation can prevent server growth from ``moving too fast, and breaking things". 

One method of regulation may be the three pillars. Deployed in the US,  Europe, and partially in Canda, for regulating and alloting tax credits to green hydrogen Cite [Federal Register, 2025; EUR-lex, 2023; Bartlett, 2024]. There's international academic literature already exploring why each of the three pillars is needed to control emissions of electrolytic hydrogen Cite [Ricks, 2023; Zeyen, 2024]. Server grid siting Cite [Vabson, 2026], regulation Cite [Park, 2026; Senga, 2026], emissions Cite [Johnson, 2026] and derived health impacts have also been explored in academic literature Cite [Mueller, 2026]. Moreover, regulators have already mobilized to regulate server loads. E3 notes dozens of large load tariffs in response to server loads weighting down the public grid Cite [Beer, 2026], the Ratepayer Protection Act passed the U.S. House of Representatives this September Cite [Howland, 2026 (2)], and various industry, and regulatory, actors are currently exploring how to manage these new sizeable data center loads Cite [Viswanath, 2026].

Our contribution is evaluating a number of regulatory policies for the regulation of new server loads. To the best of our knowledge, we are the first work to evaluate policy impacts on server load adoption throughout the continuous US. We do this at an hourly resolution and for 26 regions defined by ISO/RTO boundaries and NERC assessment areas. We produce national and per-region estimates for costs, emissions, capacity expansion and generation profile impacts of explored policies.


\section{Results}
Lower LMPs with renewables, less pollution deaths too. higher capex maybe. batteries help out. sensitivity analysis needed in projected demand and costs of fuel and generators. 

\section{Discussion}
Results vary by region 

\section{Methods}
Temoa is an large scale energy model. Our experiments use the CONUS 26 zone dataset published in may 2025 and available at Cite [OEO, 2026]. We represent the US as 26 zones roughly defined by NERC []. Demand is at an hourly resolution and sourced from [NREL] and projected to 2040 from [IEA AEO]. Generation, also at hourly resolution is regionally based. Each region has a collection of built and potential energy sources based on [FILL]. Transmisison between regions is lossy and based on the overall energy flow between these regions in the existing grid [CHECK]. The linear program behind temoa sees the demand for a representative 3360 hour (20 week) year and chooses which generation capacity to build, retire, and dispatch to fulfill the period in question. We model four periods: 2027-2030, 2030-2035, 2035-2040, and 2040-2045. The objective is minimal present cost. 

\subsection{Three Pillars in Temoa}

We build simulations including every combination of the three pillars. 

We model additionality by constraining the electricity that feeds the servers to come from renewable sources built after 2027. We create a special "server electricity" that can only be made by these eligible plants. When relaxing additionality, we do not take a plant's creation year into account. 

Hourly-matching is native to temoa; this hourly-resolution energy model enforces that electric demand meets supply at each hour. When we code the server load such that it can only be met by renewables, hourly-matching is automatically enforced. When relaxing hourly-matching, we allow servers to be fed with both server-electricity and regular electricity. If not hourly-matched, we require that the total, annual generated server-electricity from eligible plants equal the total, annual demand of the servers. 

We enforce deliverability by forcing the server load to come from generators in its same NERC [CHECK] region. We do not build transmission for server-electricity. When relaxing deliverability, we build server-transmission lines in the model such that eligible server-electricity to feed data centers can move across regions. 

We also include a server-electricity to electricity converter. This allow leftover server-electricity to flow into the regular grid and meet the load of other sectors like transportation and industry. 

One noteworthy decision here is relaxing hourly matching whenever deliverability is relaxed. This is because without deliverability, energy procurement becomes an accounting exercise rather than a simulation of the grid graph's flows. Hourly matching without deliverability would thus enforce deploying server-electricity across hours with batteries and overbuilt renewables, and doing so up to thousands of miles away from where the load is. 

\subsubsection{Partial Coverage}
We model the complete regulation of servers under the three pillars, and we also model regulating various percetages server load at every region and hour. In some cases, it is not necessary to regulate every MWh of server load to keep average costs, marginal costs, and CO2 emissions "unimpacted" from the addition of servers. 


\section{Conclusions}

Hourly matching has the largest impact in the grid's total present cost and emissions and how much generation gets built. Deliverability is also impactful, though the model herein does not capture all transmission congestion issues facing our current and future grid. Additionality, here applied as a rule that energy must come from newly-built generators, does not change total costs, emissions or built generators much. This is because within the eligible generators, there are only a few hundred [CHECK] GWh's available as of the present grid. Whether we enforce new construction of renewable generators or not, they must be built to feed the TWh's of data center load seen by the grid.


%
% Each of the commands below will create an unnumbered section with the appropriate heading.
% Remove any sections that are not relevant for your article.
% All sections except suppdata will be removed if the [anonymous] option is used.
% See iopjournal-guidelines.pdf for more information.
%

\ack{Sample text inserted for demonstration.}

\funding{Sample text inserted for demonstration.}
% This section is a list of funder names and grant numbers

\roles{Sample text inserted for demonstration.}
% List author names and the contributions made to the article, using terms from the NISO Contributor Roles Taxonomy (CRediT) https://credit.niso.org

\data{Sample text inserted for demonstration.}
% For more information on IOP Publishing's research data policy see: https://publishingsupport.iopscience.iop.org/questions/research-data/

This energy model of the contiguous United States is available at \url{https://datadocs.openenergyoutlook.org/en/latest/datasets.html}.

The code underlying our implementation of three-pillars and its analysis is available at \url{https://github.com/DavidLikesLearning/temoa-dataCenter-3Pillars.}

\suppdata{Sample text inserted for demonstration.}


\section*{References}



\end{document}


